% year: תשפ"ה
% type: בוחן
% semester: א
% moed: לדוגמה
% official_solution: yes
% source: exams/מבחנים/תשפו/midterm_example 2025.pdf
% part: א
% number: 2
% points: 20
% categories: continuity, func-limits

\begin{question}
יהיו $a,b\in\R$. נגדיר פונקציה על ידי
\[
f(x)=\begin{cases}
\frac{\sin(8x)}{x} & x<0,\\
ax+b & 0\le x\le 4,\\
\frac{x^2-3x-4}{2-\sqrt{x}} & x>4
\end{cases}
\]
לאילו ערכים של $a,b$ הפונקציה רציפה ב-$\R$ ? נמקו.
\end{question}

\begin{hint}
בכל נקודה $x\neq0,4$ הפונקציה רציפה בלי קשר ל-$a,b$. בדקו רק את נקודות התפר $x=0$ ו-$x=4$.
\end{hint}

\begin{hint}
ב-$x=4$: $x^2-3x-4=(x-4)(x+1)$, וכפלו וחלקו ב-$2+\sqrt x$.
\end{hint}

\begin{solution}
\textbf{נקודות שאינן נקודות תפר.} בכל נקודה $x<0$ הפונקציה היא מנה של פונקציות רציפות עם מכנה שונה מ-0; בכל נקודה $0<x<4$ היא פולינום; ובכל נקודה $x>4$ היא מנה של פונקציות רציפות עם מכנה $2-\sqrt x\neq0$ (כי $\sqrt x>2$). מכיוון שבכל אחת מהנקודות האלה $f$ מתלכדת עם הנוסחה בסביבה שלמה, $f$ רציפה בהן לכל $a,b$. נותר לבדוק את $x=0$ ואת $x=4$.

\textbf{הנקודה $x=0$.} $f(0)=b$. משמאל, לפי $\lim_{y\to0}\frac{\sin y}{y}=1$ עם $y=8x$:
\[
\lim_{x\to0^-}\frac{\sin(8x)}{x}=8\lim_{x\to0^-}\frac{\sin(8x)}{8x}=8\lim_{y\to0^-}\frac{\sin y}{y}=8\cdot1=8 .
\]
מימין: $\lim_{x\to0^+}(ax+b)=b=f(0)$ (רציפות פולינום). לכן $f$ רציפה ב-0 אם ורק אם $\boxed{b=8}$.

\textbf{הנקודה $x=4$.} $f(4)=4a+b=4a+8$, ומשמאל $\lim_{x\to4^-}(ax+8)=4a+8=f(4)$. מימין, גבול מהצורה $\frac00$. נכפול בצמוד:
\[
\frac{x^2-3x-4}{2-\sqrt x}\cdot\frac{2+\sqrt x}{2+\sqrt x}=\frac{(x-4)(x+1)(2+\sqrt x)}{4-x}=-(x+1)(2+\sqrt x)\qquad(x\neq4),
\]
ולכן, מרציפות הביטוי האחרון ב-4,
\[
\lim_{x\to4^+}f(x)=-(4+1)(2+\sqrt4)=-5\cdot4=-20 .
\]
$f$ רציפה ב-4 אם ורק אם $4a+8=-20$, כלומר $\boxed{a=-7}$.

\textbf{תשובה:} $f$ רציפה ב-$\R$ אם ורק אם $a=-7,\ b=8$.

(הערה: בפתרון בכתב יד המצורף הסימן "$-$" הושמט בשורה $-\lim(x+1)(2+\sqrt x)=5\cdot4=20$, אבל בשורה הבאה נכתב נכון $4a+8=-20$ והתשובה $a=-7$ נכונה.)
\end{solution}
